% handout-common.tex —— 五本讲义共享的导言区与排版宏
% 所有 handout-*.tex 通过 % handout-common.tex —— 五本讲义共享的导言区与排版宏
% 所有 handout-*.tex 通过 % handout-common.tex —— 五本讲义共享的导言区与排版宏
% 所有 handout-*.tex 通过 % handout-common.tex —— 五本讲义共享的导言区与排版宏
% 所有 handout-*.tex 通过 \input{handout-common} 引入本文件。
% 【重要】此文件为统一规范，任何单本讲义不得修改它。
% ===========================================================================
\usepackage[a4paper,left=18mm,right=18mm,top=18mm,bottom=18mm,
  includeheadfoot,headheight=15pt,headsep=4mm,footskip=10mm]{geometry}
\usepackage{amsmath,amssymb,booktabs,longtable,array,enumitem,listings,xcolor,hyperref,fancyhdr,graphicx}
\graphicspath{{figures/}}
\hypersetup{hidelinks,pdfauthor={CTF 培训资料}}

\setlength{\parindent}{2em}
\setlength{\parskip}{0.25em}
\emergencystretch=3em
\setlist{itemsep=0.15em,topsep=0.25em}
\lstset{basicstyle=\ttfamily\small,breaklines=true,columns=fullflexible,keepspaces=true,
  showstringspaces=false,frame=single,framerule=0.3pt,xleftmargin=0.4em,xrightmargin=0.4em}

% ---- 页眉页脚 -------------------------------------------------------------
\pagestyle{fancy}
\fancyhf{}
\fancyhead[L]{CTF 网络安全入门}
\fancyhead[R]{\thehandoutmark}
\fancyfoot[C]{\thepage}
\renewcommand{\headrulewidth}{0.3pt}
\newcommand{\thehandoutmark}{讲义}
% 用法（写在导言区）：\handoutmark{Web 专攻版}
\newcommand{\handoutmark}[1]{\renewcommand{\thehandoutmark}{#1}}
\newcommand{\booktitle}[1]{\hypersetup{pdftitle={#1}}}

% ---- 正文辅助宏 -----------------------------------------------------------
\newcommand{\goal}[1]{\noindent\textbf{本章目标：}#1\par}
\newcommand{\code}[1]{\texttt{#1}}
\newcommand{\kw}[1]{\textbf{#1}}
\newcommand{\flagbox}[1]{\par\noindent{\setlength{\fboxsep}{4pt}\fbox{\textbf{flag：}\texttt{#1}}}\par}
% 题面卡中的字段行：\field{附件}{...}
\newcommand{\field}[2]{\par\noindent\textbf{#1：}#2\par}

% ---- 信息框（一句话记忆 / 常见误区 / 排错指南 / 小技巧）--------------------
\newsavebox{\ibox}
\newenvironment{infobox}[1]
  {\par\medskip\begin{lrbox}{\ibox}\begin{minipage}{0.95\linewidth}\textbf{#1}\par\smallskip\hrule\smallskip}
  {\end{minipage}\end{lrbox}\noindent\fbox{\usebox{\ibox}}\par\medskip}
\newcommand{\memory}[1]{\begin{infobox}{一句话记忆}#1\end{infobox}}
\newcommand{\pitfall}[1]{\begin{infobox}{常见误区}#1\end{infobox}}
\newcommand{\debugtip}[1]{\begin{infobox}{排错指南}#1\end{infobox}}
\newcommand{\tip}[1]{\begin{infobox}{小技巧}#1\end{infobox}}

% ---- 题面卡 ---------------------------------------------------------------
% \begin{challenge}{题 R1：一句话标题} ... \field{...}{...} ... \end{challenge}
\newenvironment{challenge}[1]
  {\par\medskip\begin{lrbox}{\ibox}\begin{minipage}{0.95\linewidth}{\large\bfseries #1}\par\smallskip\hrule\smallskip}
  {\end{minipage}\end{lrbox}\noindent\fbox{\usebox{\ibox}}\par\medskip}

% ---- 插图 -----------------------------------------------------------------
% \shot[宽度比例]{文件名}{图注}{标签}   引用：\figref{标签}
\newcommand{\shot}[4][0.96]{%
  \begin{figure}[htbp]\centering
  \includegraphics[width=#1\linewidth]{#2}%
  \caption{#3}\label{fig:#4}\end{figure}}
\newcommand{\figref}[1]{图~\ref{fig:#1}}

\endinput
 引入本文件。
% 【重要】此文件为统一规范，任何单本讲义不得修改它。
% ===========================================================================
\usepackage[a4paper,left=18mm,right=18mm,top=18mm,bottom=18mm,
  includeheadfoot,headheight=15pt,headsep=4mm,footskip=10mm]{geometry}
\usepackage{amsmath,amssymb,booktabs,longtable,array,enumitem,listings,xcolor,hyperref,fancyhdr,graphicx}
\graphicspath{{figures/}}
\hypersetup{hidelinks,pdfauthor={CTF 培训资料}}

\setlength{\parindent}{2em}
\setlength{\parskip}{0.25em}
\emergencystretch=3em
\setlist{itemsep=0.15em,topsep=0.25em}
\lstset{basicstyle=\ttfamily\small,breaklines=true,columns=fullflexible,keepspaces=true,
  showstringspaces=false,frame=single,framerule=0.3pt,xleftmargin=0.4em,xrightmargin=0.4em}

% ---- 页眉页脚 -------------------------------------------------------------
\pagestyle{fancy}
\fancyhf{}
\fancyhead[L]{CTF 网络安全入门}
\fancyhead[R]{\thehandoutmark}
\fancyfoot[C]{\thepage}
\renewcommand{\headrulewidth}{0.3pt}
\newcommand{\thehandoutmark}{讲义}
% 用法（写在导言区）：\handoutmark{Web 专攻版}
\newcommand{\handoutmark}[1]{\renewcommand{\thehandoutmark}{#1}}
\newcommand{\booktitle}[1]{\hypersetup{pdftitle={#1}}}

% ---- 正文辅助宏 -----------------------------------------------------------
\newcommand{\goal}[1]{\noindent\textbf{本章目标：}#1\par}
\newcommand{\code}[1]{\texttt{#1}}
\newcommand{\kw}[1]{\textbf{#1}}
\newcommand{\flagbox}[1]{\par\noindent{\setlength{\fboxsep}{4pt}\fbox{\textbf{flag：}\texttt{#1}}}\par}
% 题面卡中的字段行：\field{附件}{...}
\newcommand{\field}[2]{\par\noindent\textbf{#1：}#2\par}

% ---- 信息框（一句话记忆 / 常见误区 / 排错指南 / 小技巧）--------------------
\newsavebox{\ibox}
\newenvironment{infobox}[1]
  {\par\medskip\begin{lrbox}{\ibox}\begin{minipage}{0.95\linewidth}\textbf{#1}\par\smallskip\hrule\smallskip}
  {\end{minipage}\end{lrbox}\noindent\fbox{\usebox{\ibox}}\par\medskip}
\newcommand{\memory}[1]{\begin{infobox}{一句话记忆}#1\end{infobox}}
\newcommand{\pitfall}[1]{\begin{infobox}{常见误区}#1\end{infobox}}
\newcommand{\debugtip}[1]{\begin{infobox}{排错指南}#1\end{infobox}}
\newcommand{\tip}[1]{\begin{infobox}{小技巧}#1\end{infobox}}

% ---- 题面卡 ---------------------------------------------------------------
% \begin{challenge}{题 R1：一句话标题} ... \field{...}{...} ... \end{challenge}
\newenvironment{challenge}[1]
  {\par\medskip\begin{lrbox}{\ibox}\begin{minipage}{0.95\linewidth}{\large\bfseries #1}\par\smallskip\hrule\smallskip}
  {\end{minipage}\end{lrbox}\noindent\fbox{\usebox{\ibox}}\par\medskip}

% ---- 插图 -----------------------------------------------------------------
% \shot[宽度比例]{文件名}{图注}{标签}   引用：\figref{标签}
\newcommand{\shot}[4][0.96]{%
  \begin{figure}[htbp]\centering
  \includegraphics[width=#1\linewidth]{#2}%
  \caption{#3}\label{fig:#4}\end{figure}}
\newcommand{\figref}[1]{图~\ref{fig:#1}}

\endinput
 引入本文件。
% 【重要】此文件为统一规范，任何单本讲义不得修改它。
% ===========================================================================
\usepackage[a4paper,left=18mm,right=18mm,top=18mm,bottom=18mm,
  includeheadfoot,headheight=15pt,headsep=4mm,footskip=10mm]{geometry}
\usepackage{amsmath,amssymb,booktabs,longtable,array,enumitem,listings,xcolor,hyperref,fancyhdr,graphicx}
\graphicspath{{figures/}}
\hypersetup{hidelinks,pdfauthor={CTF 培训资料}}

\setlength{\parindent}{2em}
\setlength{\parskip}{0.25em}
\emergencystretch=3em
\setlist{itemsep=0.15em,topsep=0.25em}
\lstset{basicstyle=\ttfamily\small,breaklines=true,columns=fullflexible,keepspaces=true,
  showstringspaces=false,frame=single,framerule=0.3pt,xleftmargin=0.4em,xrightmargin=0.4em}

% ---- 页眉页脚 -------------------------------------------------------------
\pagestyle{fancy}
\fancyhf{}
\fancyhead[L]{CTF 网络安全入门}
\fancyhead[R]{\thehandoutmark}
\fancyfoot[C]{\thepage}
\renewcommand{\headrulewidth}{0.3pt}
\newcommand{\thehandoutmark}{讲义}
% 用法（写在导言区）：\handoutmark{Web 专攻版}
\newcommand{\handoutmark}[1]{\renewcommand{\thehandoutmark}{#1}}
\newcommand{\booktitle}[1]{\hypersetup{pdftitle={#1}}}

% ---- 正文辅助宏 -----------------------------------------------------------
\newcommand{\goal}[1]{\noindent\textbf{本章目标：}#1\par}
\newcommand{\code}[1]{\texttt{#1}}
\newcommand{\kw}[1]{\textbf{#1}}
\newcommand{\flagbox}[1]{\par\noindent{\setlength{\fboxsep}{4pt}\fbox{\textbf{flag：}\texttt{#1}}}\par}
% 题面卡中的字段行：\field{附件}{...}
\newcommand{\field}[2]{\par\noindent\textbf{#1：}#2\par}

% ---- 信息框（一句话记忆 / 常见误区 / 排错指南 / 小技巧）--------------------
\newsavebox{\ibox}
\newenvironment{infobox}[1]
  {\par\medskip\begin{lrbox}{\ibox}\begin{minipage}{0.95\linewidth}\textbf{#1}\par\smallskip\hrule\smallskip}
  {\end{minipage}\end{lrbox}\noindent\fbox{\usebox{\ibox}}\par\medskip}
\newcommand{\memory}[1]{\begin{infobox}{一句话记忆}#1\end{infobox}}
\newcommand{\pitfall}[1]{\begin{infobox}{常见误区}#1\end{infobox}}
\newcommand{\debugtip}[1]{\begin{infobox}{排错指南}#1\end{infobox}}
\newcommand{\tip}[1]{\begin{infobox}{小技巧}#1\end{infobox}}

% ---- 题面卡 ---------------------------------------------------------------
% \begin{challenge}{题 R1：一句话标题} ... \field{...}{...} ... \end{challenge}
\newenvironment{challenge}[1]
  {\par\medskip\begin{lrbox}{\ibox}\begin{minipage}{0.95\linewidth}{\large\bfseries #1}\par\smallskip\hrule\smallskip}
  {\end{minipage}\end{lrbox}\noindent\fbox{\usebox{\ibox}}\par\medskip}

% ---- 插图 -----------------------------------------------------------------
% \shot[宽度比例]{文件名}{图注}{标签}   引用：\figref{标签}
\newcommand{\shot}[4][0.96]{%
  \begin{figure}[htbp]\centering
  \includegraphics[width=#1\linewidth]{#2}%
  \caption{#3}\label{fig:#4}\end{figure}}
\newcommand{\figref}[1]{图~\ref{fig:#1}}

\endinput
 引入本文件。
% 【重要】此文件为统一规范，任何单本讲义不得修改它。
% ===========================================================================
\usepackage[a4paper,left=18mm,right=18mm,top=18mm,bottom=18mm,
  includeheadfoot,headheight=15pt,headsep=4mm,footskip=10mm]{geometry}
\usepackage{amsmath,amssymb,booktabs,longtable,array,enumitem,listings,xcolor,hyperref,fancyhdr,graphicx}
\graphicspath{{figures/}}
\hypersetup{hidelinks,pdfauthor={CTF 培训资料}}

\setlength{\parindent}{2em}
\setlength{\parskip}{0.25em}
\emergencystretch=3em
\setlist{itemsep=0.15em,topsep=0.25em}
\lstset{basicstyle=\ttfamily\small,breaklines=true,columns=fullflexible,keepspaces=true,
  showstringspaces=false,frame=single,framerule=0.3pt,xleftmargin=0.4em,xrightmargin=0.4em}

% ---- 页眉页脚 -------------------------------------------------------------
\pagestyle{fancy}
\fancyhf{}
\fancyhead[L]{CTF 网络安全入门}
\fancyhead[R]{\thehandoutmark}
\fancyfoot[C]{\thepage}
\renewcommand{\headrulewidth}{0.3pt}
\newcommand{\thehandoutmark}{讲义}
% 用法（写在导言区）：\handoutmark{Web 专攻版}
\newcommand{\handoutmark}[1]{\renewcommand{\thehandoutmark}{#1}}
\newcommand{\booktitle}[1]{\hypersetup{pdftitle={#1}}}

% ---- 正文辅助宏 -----------------------------------------------------------
\newcommand{\goal}[1]{\noindent\textbf{本章目标：}#1\par}
\newcommand{\code}[1]{\texttt{#1}}
\newcommand{\kw}[1]{\textbf{#1}}
\newcommand{\flagbox}[1]{\par\noindent{\setlength{\fboxsep}{4pt}\fbox{\textbf{flag：}\texttt{#1}}}\par}
% 题面卡中的字段行：\field{附件}{...}
\newcommand{\field}[2]{\par\noindent\textbf{#1：}#2\par}

% ---- 信息框（一句话记忆 / 常见误区 / 排错指南 / 小技巧）--------------------
\newsavebox{\ibox}
\newenvironment{infobox}[1]
  {\par\medskip\begin{lrbox}{\ibox}\begin{minipage}{0.95\linewidth}\textbf{#1}\par\smallskip\hrule\smallskip}
  {\end{minipage}\end{lrbox}\noindent\fbox{\usebox{\ibox}}\par\medskip}
\newcommand{\memory}[1]{\begin{infobox}{一句话记忆}#1\end{infobox}}
\newcommand{\pitfall}[1]{\begin{infobox}{常见误区}#1\end{infobox}}
\newcommand{\debugtip}[1]{\begin{infobox}{排错指南}#1\end{infobox}}
\newcommand{\tip}[1]{\begin{infobox}{小技巧}#1\end{infobox}}

% ---- 题面卡 ---------------------------------------------------------------
% \begin{challenge}{题 R1：一句话标题} ... \field{...}{...} ... \end{challenge}
\newenvironment{challenge}[1]
  {\par\medskip\begin{lrbox}{\ibox}\begin{minipage}{0.95\linewidth}{\large\bfseries #1}\par\smallskip\hrule\smallskip}
  {\end{minipage}\end{lrbox}\noindent\fbox{\usebox{\ibox}}\par\medskip}

% ---- 插图 -----------------------------------------------------------------
% \shot[宽度比例]{文件名}{图注}{标签}   引用：\figref{标签}
\newcommand{\shot}[4][0.96]{%
  \begin{figure}[htbp]\centering
  \includegraphics[width=#1\linewidth]{#2}%
  \caption{#3}\label{fig:#4}\end{figure}}
\newcommand{\figref}[1]{图~\ref{fig:#1}}

\endinput
